\section{Conclusion}
Verified fallback-free placement of a modern detector on DLA is achievable, and
at the measured operating point it costs more in the application pipeline than
the same network on GPU, with the boundary rather than the accelerator
accounting for the difference.
For YOLO11, the resulting adaptation combines C2DLA with a static raw prediction boundary;
fallback-disabled compilation and engine inspection verify one DLA loadable
in FP16 and INT8 on the tested Orin/TensorRT stack. FP16 nearly preserves
adapted GPU aggregate AP. Relative to FP16, INT8 reduces strict DLA engine-plus-transfer latency,
application-pipeline latency, and module energy at 10\,Hz, with an accuracy cost.
Adapted GPU remains faster and lower-energy at this operating point.

The case study shows that DLA adaptation is not only an operator-replacement
problem. The prediction boundary, tensor layouts, host processing, and placement
checks also matter. Concurrent robotic workloads remain outside this evaluation. Whether removing
the learned detector from GPU compute benefits them is the open question this
case study raises but does not answer: the measured 7--10\% GR3D utilization
shows the GPU is not idle during strict execution, and the strict configuration
occupies the platform for longer per frame, so a system-level test with a
competing workload is required before any such benefit can be claimed.
Other models and DLA configurations require their own validation.
